\documentclass{article}

\usepackage[margin=1in]{geometry}
\usepackage{amsmath,amssymb}
\usepackage{braket}
\usepackage{booktabs}
\usepackage{hyperref}

% Code formatting
\usepackage{xcolor}
\usepackage{listings}
\lstdefinelanguage{Mathematica}{
  morekeywords={Module,Table,If,Do,Join,AppendTo,Length,ConstantArray,
    RowReduce,Simplify,Expand,Reduce,Minimize,Maximize,Complement,
    Range,Position,Total,Sum,And,Or,Symbol,Association,Print,
    StringLength,StringPart,True,False},
  sensitive=true,
  morecomment=[s]{(*}{*)},
  morestring=[b]"
}
\lstset{
  language=Mathematica,
  basicstyle=\ttfamily\small,
  keywordstyle=\color{blue!60!black},
  commentstyle=\color{green!40!black},
  stringstyle=\color{red!50!black},
  showstringspaces=false,
  breaklines=true,
  frame=single,
  columns=fullflexible
}

% shorthand
\newcommand{\SigZ}{\Sigma_Z}
\newcommand{\conv}{\mathrm{conv}}

\title{Independent Wolfram Language Verification \\
  of KL-Derived Linear Systems from \texttt{LPv3.tex}}
\author{}
\date{}

\begin{document}
\maketitle
\tableofcontents

%%%%%%%%%%%%%%%%%%%%%%%%%%%%%%%%%%%%%%%%%%%%%%%%%%%%%%%%%%%%%%%%
\section{Introduction}
%%%%%%%%%%%%%%%%%%%%%%%%%%%%%%%%%%%%%%%%%%%%%%%%%%%%%%%%%%%%%%%%

The companion note \texttt{LPv3.tex} derives parametric solutions of the Knill--Laflamme (KL) $Z$-marginal conditions for six quantum-code examples (five one-parameter families A--E and a worked $((5,2,2))$ example) using Python/SymPy.
This document reports an independent, end-to-end re-derivation and verification of \emph{every} numerical and algebraic claim in \texttt{LPv3.tex}, carried out entirely in Wolfram Language (Mathematica).
The accompanying script \texttt{verify\_LP.wl} (425~lines) is self-contained and can be run via
\begin{center}
\texttt{wolframscript -file verify\_LP.wl}
\end{center}

\medskip\noindent
\textbf{Result.} \emph{Zero discrepancies.}
All parametric solutions, feasibility regions, $Z$-marginal vectors, $\SigZ$ polynomials, $\SigZ$ ranges, and convex-hull dimension tests match exactly.

\subsection{Scope of verification}

The verification pipeline, applied to each of the six examples, checks:
\begin{enumerate}
  \item Build the sign matrix $s_i(x)=(-1)^{x_i}$ for each codeword.
  \item Construct the augmented matrix $[A\mid b]$ encoding $Z$-marginal equalities and normalization.
  \item Compute \texttt{RowReduce} (RREF over $\mathbb{Q}$) to get the affine parametric solution.
  \item Identify pivot/free columns, extract the parametric solution.
  \item Impose nonnegativity via \texttt{Reduce} to obtain the feasibility region.
  \item Compute $Z$-marginal vectors from \emph{both} $C_0$ and $C_1$ sides and confirm they match.
  \item Compute $\SigZ=\sum_i \langle Z_i\rangle^2$ and its range via \texttt{Minimize}/\texttt{Maximize}.
  \item (Appendix A.3 test) Minimize and maximize each $\langle Z_i\rangle$ to determine whether $\conv(V_0)\cap\conv(V_1)$ is a point or has positive dimension.
\end{enumerate}

\subsection{Relationship to the Lean formalization}

Families A--E and the $(5,2,2)$ example are also formalized in Lean~4 (Mathlib).
The Lean proofs verify that the \emph{given} probability assignments satisfy the SS condition, residue-shift screen, normalization, and $Z$-type KL equalities---that is, \textbf{correctness}: ``these formulas work.''
The Wolfram verification additionally confirms \textbf{completeness}: ``these are \emph{all} solutions,'' by deriving the parametric families from scratch via RREF and showing the feasibility regions are exactly as claimed.

\begin{center}
\begin{tabular}{lcc}
\toprule
Property & Lean & Wolfram \\
\midrule
SS support condition & \checkmark & --- \\
Residue-shift screen & \checkmark & --- \\
Probability normalization & \checkmark & \checkmark \\
$Z$-type KL ($\langle Z_i\rangle$ matches target) & \checkmark & \checkmark \\
$\SigZ$ range & \checkmark & \checkmark \\
Parametric solution is complete (RREF) & --- & \checkmark \\
Feasibility region is exact & --- & \checkmark \\
$Z$-marginals from $C_0$ and $C_1$ agree & --- & \checkmark \\
$\conv(V_0)\cap\conv(V_1)$ dimension & --- & \checkmark \\
\bottomrule
\end{tabular}
\end{center}

%%%%%%%%%%%%%%%%%%%%%%%%%%%%%%%%%%%%%%%%%%%%%%%%%%%%%%%%%%%%%%%%
\section{Summary of results}
\label{sec:summary}
%%%%%%%%%%%%%%%%%%%%%%%%%%%%%%%%%%%%%%%%%%%%%%%%%%%%%%%%%%%%%%%%

Table~\ref{tab:summary} gives the headline numbers. Full per-family comparisons follow in Section~\ref{sec:families}.

\begin{table}[ht]
\centering
\begin{tabular}{lcccccl}
\toprule
Example & Vars & Rank & Free & Feasibility & $\SigZ$ range & $\conv$ test \\
\midrule
Family A & 6 & 5 & 1 & $0\le t\le\tfrac12$ & $[\tfrac12,1]$ & segment \\
Family B & 8 & 7 & 1 & $0\le t\le\tfrac13$ & $[\tfrac13,\tfrac59]$ & segment \\
Family C & 8 & 7 & 1 & $0\le t\le\tfrac13$ & $[\tfrac13,\tfrac59]$ & segment \\
Family D & 8 & 7 & 1 & $0\le t\le\tfrac13$ & $[\tfrac13,\tfrac59]$ & segment \\
Family E & 6 & 5 & 1 & $0\le t\le\tfrac12$ & $[\tfrac12,1]$ & segment \\
Worked $(5,2,2)$ & 9 & 7 & 2 & triangle & $\{\tfrac37\}$ & point \\
\bottomrule
\end{tabular}
\caption{Summary of verification results. ``Vars'' is $|C_0|+|C_1|$; ``Rank'' is the rank of the augmented system after RREF; ``Free'' is the number of free parameters. For the worked example, the feasibility region is $u\ge0, v\ge0, u+v\le\tfrac{2}{7}$. All entries match \texttt{LPv3.tex}.}
\label{tab:summary}
\end{table}

%%%%%%%%%%%%%%%%%%%%%%%%%%%%%%%%%%%%%%%%%%%%%%%%%%%%%%%%%%%%%%%%
\section{Family-by-family verification}
\label{sec:families}
%%%%%%%%%%%%%%%%%%%%%%%%%%%%%%%%%%%%%%%%%%%%%%%%%%%%%%%%%%%%%%%%

For each family we list the supports, the parametric solution obtained from \texttt{RowReduce}, the feasibility region from \texttt{Reduce}, the $Z$-marginal vector, $\SigZ$ and its range, and the comparison with \texttt{LPv3.tex}.
In every case the results agree exactly.

%% ---- Family A ----
\subsection{Family A: $(m,\mathbf{w},S_0,S_1)=(6,(4,5,1,2,2),0,3)$}

\paragraph{Supports.}
$C_0=\{10010, 11101, 01100\}$,\quad $C_1=\{11000, 10111, 00110\}$.

\paragraph{Linear system.}
6 variables, rank~5, 1 free parameter.
Free variable: $t:=Y_{00110}$.

\paragraph{Parametric solution (Wolfram output).}
\[
\begin{aligned}
X_{10010} &= \tfrac{1}{2}, & X_{11101} &= \tfrac{1}{2}-t, & X_{01100} &= t, \\
Y_{11000} &= \tfrac{1}{2}, & Y_{10111} &= \tfrac{1}{2}-t, & Y_{00110} &= t.
\end{aligned}
\]

\paragraph{Feasibility.} $0\le t\le \frac{1}{2}$. \checkmark

\paragraph{$Z$-marginals.}
Computed independently from both $C_0$ and $C_1$:
\[
(\langle Z_1\rangle,\dots,\langle Z_5\rangle)
= (2t-1, 0, 0, 0, 2t). \quad\checkmark
\]

\paragraph{$\SigZ$.}
$\SigZ(t)=8t^2-4t+1$, range $[\tfrac12,1]$, minimum at $t=\tfrac14$. \checkmark

%% ---- Family B ----
\subsection{Family B: $(m,\mathbf{w},S_0,S_1)=(6,(1,1,1,5,3),0,2)$}

\paragraph{Supports.}
$C_0=\{00000, 00110, 10010, 11101\}$,\quad
$C_1=\{00011, 01100, 10100, 11000\}$.

\paragraph{Linear system.}
8 variables, rank~7, 1 free parameter.
Free variable: $t:=Y_{11000}$.

\paragraph{Parametric solution.}
\[
\begin{aligned}
X_{00000}=X_{11101}=Y_{00011}=Y_{10100} &= \tfrac{1}{3},\\
X_{00110}=Y_{01100} &= \tfrac{1}{3}-t,\\
X_{10010}=Y_{11000} &= t.
\end{aligned}
\]

\paragraph{Feasibility.} $0\le t\le \frac{1}{3}$. \checkmark

\paragraph{$Z$-marginals.}
\[
(\langle Z_1\rangle,\dots,\langle Z_5\rangle)
= (\tfrac13-2t, \tfrac13, {-}\tfrac13+2t, \tfrac13, \tfrac13). \quad\checkmark
\]

\paragraph{$\SigZ$.}
$\SigZ(t)=8t^2-\frac{8}{3}t+\frac{5}{9}$, range $[\tfrac13,\tfrac59]$, minimum at $t=\tfrac16$. \checkmark

%% ---- Family C ----
\subsection{Family C: $(m,\mathbf{w},S_0,S_1)=(6,(3,3,1,1,1),0,4)$}

\paragraph{Supports.}
$C_0=\{00000, 01111, 10111, 11000\}$,\quad
$C_1=\{01001, 01100, 10010, 10100\}$.

\paragraph{Linear system.}
8 variables, rank~7, 1 free parameter.
Free variable: $t:=Y_{10100}$.

\paragraph{Parametric solution.}
\[
\begin{aligned}
X_{00000}=X_{11000}=Y_{01001}=Y_{10010} &= \tfrac{1}{3},\\
X_{01111}=Y_{01100} &= \tfrac{1}{3}-t,\\
X_{10111}=Y_{10100} &= t.
\end{aligned}
\]

\paragraph{Feasibility.} $0\le t\le \frac{1}{3}$. \checkmark

\paragraph{$Z$-marginals.}
\[
(\langle Z_1\rangle,\dots,\langle Z_5\rangle)
= (\tfrac13-2t, {-}\tfrac13+2t, \tfrac13, \tfrac13, \tfrac13). \quad\checkmark
\]

\paragraph{$\SigZ$.}
$\SigZ(t)=8t^2-\frac{8}{3}t+\frac{5}{9}$, range $[\tfrac13,\tfrac59]$, minimum at $t=\tfrac16$. \checkmark

%% ---- Family D ----
\subsection{Family D: $(m,\mathbf{w},S_0,S_1)=(6,(3,1,3,5,5),0,4)$}

\paragraph{Supports.}
$C_0=\{00000, 01001, 11101, 11110\}$,\quad
$C_1=\{00011, 01100, 10111, 11000\}$.

\paragraph{Linear system.}
8 variables, rank~7, 1 free parameter.
Free variable: $t:=Y_{10111}$.

\paragraph{Parametric solution.}
\[
\begin{aligned}
X_{00000}=X_{11110}=Y_{01100}=Y_{11000} &= \tfrac{1}{3},\\
X_{01001}=Y_{00011} &= \tfrac{1}{3}-t,\\
X_{11101}=Y_{10111} &= t.
\end{aligned}
\]

\paragraph{Feasibility.} $0\le t\le \frac{1}{3}$. \checkmark

\paragraph{$Z$-marginals.}
\[
(\langle Z_1\rangle,\dots,\langle Z_5\rangle)
= (\tfrac13-2t, {-}\tfrac13, \tfrac13-2t, \tfrac13, \tfrac13). \quad\checkmark
\]

\paragraph{$\SigZ$.}
$\SigZ(t)=8t^2-\frac{8}{3}t+\frac{5}{9}$, range $[\tfrac13,\tfrac59]$, minimum at $t=\tfrac16$. \checkmark

%% ---- Family E ----
\subsection{Family E: $(m,\mathbf{w},S_0,S_1)=(8,(1,6,2,5,5),0,4)$}

\paragraph{Supports.}
$C_0=\{01011, 10110, 10101\}$,\quad $C_1=\{00111, 11010, 11001\}$.

\paragraph{Linear system.}
6 variables, rank~5, 1 free parameter.
Free variable: $t:=Y_{11001}$.

\paragraph{Parametric solution.}
\[
\begin{aligned}
X_{01011} = Y_{00111} &= \tfrac{1}{2}, \\
X_{10110} = Y_{11010} &= \tfrac{1}{2}-t, \\
X_{10101} = Y_{11001} &= t.
\end{aligned}
\]

\paragraph{Feasibility.} $0\le t\le \frac{1}{2}$. \checkmark

\paragraph{$Z$-marginals.}
\[
(\langle Z_1\rangle,\dots,\langle Z_5\rangle)
= (0, 0, 0, {-}1+2t, {-}2t). \quad\checkmark
\]

\paragraph{$\SigZ$.}
$\SigZ(t)=8t^2-4t+1$, range $[\tfrac12,1]$, minimum at $t=\tfrac14$. \checkmark

%% ---- Worked example ----
\subsection{Worked example: $((5,2,2))$, $m=7$}

\paragraph{Supports.}
$C_0=\{00000, 01111, 10111\}$,\quad
$C_1=\{00011, 00101, 00110, 11001, 11010, 11100\}$.

\paragraph{Linear system.}
9 variables, rank~7, 2 free parameters.
Free variables: $u:=Y_{11010}$, $v:=Y_{11100}$.

\paragraph{Parametric solution.}
\[
X_{00000}=\tfrac{3}{7},\quad X_{01111}=X_{10111}=\tfrac{2}{7},
\]
\[
\begin{aligned}
Y_{00011}&=\tfrac{1}{7}+v,\qquad
Y_{00101}=\tfrac{1}{7}+u,\qquad
Y_{00110}=\tfrac{3}{7}-u-v,\\
Y_{11001}&=\tfrac{2}{7}-u-v,\qquad
Y_{11010}=u,\qquad
Y_{11100}=v.
\end{aligned}
\]

\paragraph{Feasibility.} $u\ge 0$, $v\ge 0$, $u+v\le\frac{2}{7}$. \checkmark

\paragraph{$Z$-marginals.}
Constant (independent of $u,v$):
\[
(\langle Z_1\rangle,\dots,\langle Z_5\rangle)
= \left(\tfrac{3}{7}, \tfrac{3}{7}, {-}\tfrac{1}{7}, {-}\tfrac{1}{7}, {-}\tfrac{1}{7}\right). \quad\checkmark
\]

\paragraph{$\SigZ$.}
$\SigZ = 2(\tfrac{3}{7})^2+3(\tfrac{1}{7})^2 = \tfrac{3}{7}$ (constant). \checkmark

%%%%%%%%%%%%%%%%%%%%%%%%%%%%%%%%%%%%%%%%%%%%%%%%%%%%%%%%%%%%%%%%
\section{Convex-hull dimension test (Appendix A.3 of \texttt{LPv3.tex})}
\label{sec:convhull}
%%%%%%%%%%%%%%%%%%%%%%%%%%%%%%%%%%%%%%%%%%%%%%%%%%%%%%%%%%%%%%%%

The ``practical test for varying marginals'' from Appendix~A.3 of \texttt{LPv3.tex} is:
for each example, minimize and maximize each $Z$-marginal coordinate $\langle Z_i\rangle$ subject to all feasibility constraints (equalities from RREF plus nonnegativity).
If $\min\langle Z_i\rangle\neq\max\langle Z_i\rangle$ for any $i$, then $\conv(V_0)\cap\conv(V_1)$ has positive dimension (marginals vary); if all are equal, the intersection is a single point.

We implemented this using \texttt{Minimize}/\texttt{Maximize} on the parametric $Z$-marginal expressions subject to the \texttt{Reduce}-derived feasibility region.

\begin{table}[ht]
\centering
\begin{tabular}{lcccccl}
\toprule
Example & $\langle Z_1\rangle$ & $\langle Z_2\rangle$ & $\langle Z_3\rangle$ & $\langle Z_4\rangle$ & $\langle Z_5\rangle$ & Verdict \\
\midrule
Family A & $[-1,0]$ & $\{0\}$ & $\{0\}$ & $\{0\}$ & $[0,1]$ & \textbf{varies} \\
Family B & $[-\frac13,\frac13]$ & $\{\frac13\}$ & $[-\frac13,\frac13]$ & $\{\frac13\}$ & $\{\frac13\}$ & \textbf{varies} \\
Family C & $[-\frac13,\frac13]$ & $[-\frac13,\frac13]$ & $\{\frac13\}$ & $\{\frac13\}$ & $\{\frac13\}$ & \textbf{varies} \\
Family D & $[-\frac13,\frac13]$ & $\{-\frac13\}$ & $[-\frac13,\frac13]$ & $\{\frac13\}$ & $\{\frac13\}$ & \textbf{varies} \\
Family E & $\{0\}$ & $\{0\}$ & $\{0\}$ & $[-1,0]$ & $[-1,0]$ & \textbf{varies} \\
Worked & $\{\frac37\}$ & $\{\frac37\}$ & $\{-\frac17\}$ & $\{-\frac17\}$ & $\{-\frac17\}$ & \textbf{fixed} \\
\bottomrule
\end{tabular}
\caption{Dimension test: ranges of individual $Z$-marginal coordinates over the feasibility region. Intervals denote varying coordinates; singletons denote fixed ones. All agree with \texttt{LPv3.tex}.}
\label{tab:dimtest}
\end{table}

\paragraph{Interpretation.}
\begin{itemize}
\item \textbf{Families A--E}: $\conv(V_0)\cap\conv(V_1)$ is a line segment in $\mathbb{R}^5$, parameterized by $t$.
  Two of the five $\langle Z_i\rangle$ coordinates vary (the rest are constant);
  the varying pair always involves the coordinates whose sign vectors differ across the free-parameter codeword pair.
\item \textbf{Worked $((5,2,2))$}: $\conv(V_0)\cap\conv(V_1)$ is a single point $(\frac{3}{7},\frac{3}{7},-\frac{1}{7},-\frac{1}{7},-\frac{1}{7})$.
  The two free parameters $(u,v)$ only change how this fixed point is decomposed as a convex combination of vertices in~$V_1$; they do not move the marginal vector itself.
\end{itemize}

%%%%%%%%%%%%%%%%%%%%%%%%%%%%%%%%%%%%%%%%%%%%%%%%%%%%%%%%%%%%%%%%
\section{Structural observations}
\label{sec:observations}
%%%%%%%%%%%%%%%%%%%%%%%%%%%%%%%%%%%%%%%%%%%%%%%%%%%%%%%%%%%%%%%%

The Wolfram computation confirms the following structural patterns noted in \texttt{LPv3.tex} and \texttt{lambdav2.tex}.

\begin{enumerate}
\item \textbf{Families A and E} (3+3 supports) share the same $\SigZ$ polynomial
  $8t^2-4t+1$, with range $[\frac{1}{2},1]$ over the respective feasibility interval.

\item \textbf{Families B, C, and D} (4+4 supports) share the same $\SigZ$ polynomial
  $8t^2-\frac{8}{3}t+\frac{5}{9}$, with range $[\frac{1}{3},\frac{5}{9}]$.

\item For all five one-parameter families, the $Z$-marginal vector varies continuously with $t$, confirming that $\conv(V_0)\cap\conv(V_1)$ is a line segment (not a point).

\item For the worked $((5,2,2))$ example, the two free parameters produce a 2-simplex (triangle) in probability space, but the $Z$-marginals are constant, so $\conv(V_0)\cap\conv(V_1)$ collapses to a single point.

\item All computations are exact over $\mathbb{Q}$---no floating point is used at any stage.
\end{enumerate}

%%%%%%%%%%%%%%%%%%%%%%%%%%%%%%%%%%%%%%%%%%%%%%%%%%%%%%%%%%%%%%%%
\section{Method details}
\label{sec:method}
%%%%%%%%%%%%%%%%%%%%%%%%%%%%%%%%%%%%%%%%%%%%%%%%%%%%%%%%%%%%%%%%

\subsection{Core operation: \texttt{RowReduce}}

The central computation is \texttt{RowReduce} on the augmented matrix $[A\mid b]$ over~$\mathbb{Q}$.
The matrix has $n+2$ rows ($n$ shared-marginal equations plus 2 normalizations) and $|C_0|+|C_1|+1$ columns (one per probability variable, plus the RHS).
The $(i,j)$ entry in the marginal block is $(-1)^{(x_j)_i}$ for $C_0$-columns and $-(-1)^{(y_k)_i}$ for $C_1$-columns.

After RREF, pivot columns identify determined variables and free columns identify parameters.
The parametric solution is read off directly from the reduced rows.

\subsection{Feasibility via \texttt{Reduce}}

Substituting the parametric expressions into $p_j\ge 0$ for every probability variable yields a conjunction of linear inequalities in the free parameters.
Calling \texttt{Reduce[\ldots, freeParams, Reals]} produces a simplified description.
For the one-parameter families, this is always an interval $[0,t_{\max}]$.
For the two-parameter worked example, \texttt{Reduce} returns a case-split form; equivalence with the simpler triangle description $u\ge0$, $v\ge0$, $u+v\le\frac{2}{7}$ was confirmed using \texttt{Resolve[Exists[\ldots], Reals]}.

\subsection{$Z$-marginals and cross-check}

The $Z$-marginal vector is computed independently from both the $C_0$ and $C_1$ sides:
\[
\langle Z_i\rangle_{C_0}
= \sum_{x\in C_0} X_x\,(-1)^{x_i},
\qquad
\langle Z_i\rangle_{C_1}
= \sum_{y\in C_1} Y_y\,(-1)^{y_i}.
\]
In every example, \texttt{Simplify[zVec0 - zVec1]} returns the zero vector, confirming the shared-marginal equations are satisfied.

\subsection{Optimization}

$\SigZ$ ranges are obtained by calling \texttt{Minimize} and \texttt{Maximize} on $\sum_i \langle Z_i\rangle^2$ subject to the feasibility region.
The same approach is used for the dimension test (Section~\ref{sec:convhull}), applied to individual $\langle Z_i\rangle$ coordinates.

%%%%%%%%%%%%%%%%%%%%%%%%%%%%%%%%%%%%%%%%%%%%%%%%%%%%%%%%%%%%%%%%
\appendix
\section{Wolfram Language code listing}
\label{sec:code}
%%%%%%%%%%%%%%%%%%%%%%%%%%%%%%%%%%%%%%%%%%%%%%%%%%%%%%%%%%%%%%%%

The complete verification script is in \texttt{verify\_LP.wl} (425~lines).
Below we reproduce the core reusable functions; the full file additionally contains the per-example invocations and the dimension test from Section~\ref{sec:convhull}.

\subsection{Sign vector construction}

\begin{lstlisting}
bitToSign[s_String] := Table[
  If[StringPart[s, i] == "0", 1, -1],
  {i, StringLength[s]}
]
\end{lstlisting}

\subsection{Build augmented matrix and solve via RREF}

\begin{lstlisting}
solveKL[c0Strings_List, c1Strings_List] := Module[
  {n, n0, n1, nVars, sC0, sC1, aug, rref,
   pivotCols, freeCols, varNames, pos, p},

  n  = StringLength[c0Strings[[1]]];
  n0 = Length[c0Strings];
  n1 = Length[c1Strings];
  nVars = n0 + n1;

  sC0 = bitToSign /@ c0Strings;
  sC1 = bitToSign /@ c1Strings;

  (* Augmented matrix [A | b]:
     Rows 1..n:  Z_i shared-marginal equations
     Row n+1:    sum X_j = 1
     Row n+2:    sum Y_k = 1  *)
  aug = Table[
    Join[Table[sC0[[j, i]], {j, n0}],
         Table[-sC1[[k, i]], {k, n1}], {0}],
    {i, n}];
  AppendTo[aug,
    Join[ConstantArray[1, n0], ConstantArray[0, n1], {1}]];
  AppendTo[aug,
    Join[ConstantArray[0, n0], ConstantArray[1, n1], {1}]];

  rref = RowReduce[aug];

  (* Identify pivots *)
  pivotCols = {};
  Do[pos = Position[rref[[row]], _?(# != 0 &)];
    If[Length[pos] > 0,
      p = pos[[1, 1]];
      If[p <= nVars, AppendTo[pivotCols, p]]],
    {row, Length[rref]}];
  freeCols = Complement[Range[nVars], pivotCols];

  varNames = Join[("X_" <> # &) /@ c0Strings,
                  ("Y_" <> # &) /@ c1Strings];

  <| "n`` -> n, ''n0`` -> n0, ''n1`` -> n1,
     ''nVars`` -> nVars, ''rref`` -> rref,
     ''pivotCols`` -> pivotCols, ''freeCols" -> freeCols,
     ``varNames'' -> varNames,
     ``rank'' -> Length[pivotCols],
     ``nFree'' -> Length[freeCols],
     "C0`` -> c0Strings, ''C1" -> c1Strings,
     ``sC0'' -> sC0, ``sC1'' -> sC1 |>
]
\end{lstlisting}

\subsection{Extract parametric solution}

\begin{lstlisting}
extractSolution[result_Association] := Module[
  {rref, nVars, pivotCols, freeCols, varNames,
   freeSyms, sol, row, pc, expr},

  rref = result[``rref''];  nVars = result[``nVars''];
  pivotCols = result[``pivotCols''];
  freeCols  = result[``freeCols''];
  varNames  = result[``varNames''];

  freeSyms = Table[Symbol["t" <> ToString[k]],
                   {k, Length[freeCols]}];
  sol = <||>;

  (* Free variables *)
  Do[sol[varNames[[freeCols[[k]]]]] = freeSyms[[k]],
    {k, Length[freeCols]}];

  (* Pivot variables from RREF *)
  Do[row = rref[[r]]; pc = pivotCols[[r]];
    If[pc <= nVars,
      expr = row[[nVars + 1]];
      Do[expr = expr - row[[freeCols[[k]]]] * freeSyms[[k]],
        {k, Length[freeCols]}];
      sol[varNames[[pc]]] = Simplify[expr]],
    {r, Length[pivotCols]}];

  <| ``solution'' -> sol, ``freeSyms'' -> freeSyms,
     ``freeNames'' -> (varNames[[#]] & /@ freeCols) |>
]
\end{lstlisting}

\subsection{Compute $Z$-marginals and $\SigZ$}

\begin{lstlisting}
computeZMarginals[result_Association,
                  solData_Association] := Module[
  {n, sC0, sC1, c0, c1, sol,
   probs0, probs1, zVec0, zVec1, sigZ},

  n = result["n``]; sC0 = result[''sC0"];
  sC1 = result[``sC1''];
  c0 = result["C0``]; c1 = result[''C1"];
  sol = solData[``solution''];

  probs0 = Table[sol["X_" <> c0[[j]]], {j, Length[c0]}];
  probs1 = Table[sol["Y_" <> c1[[k]]], {k, Length[c1]}];

  zVec0 = Table[
    Sum[sC0[[j, i]] * probs0[[j]], {j, Length[c0]}],
    {i, n}] // Simplify;
  zVec1 = Table[
    Sum[sC1[[k, i]] * probs1[[k]], {k, Length[c1]}],
    {i, n}] // Simplify;

  sigZ = Total[zVec0^2] // Expand // Simplify;

  <| ``zMargC0'' -> zVec0, ``zMargC1'' -> zVec1,
     ``zMargMatch'' ->
       (Simplify[zVec0 - zVec1] ===
        ConstantArray[0, n]),
     ``sigmaZ'' -> sigZ |>
]
\end{lstlisting}

\subsection{Feasibility region}

\begin{lstlisting}
computeFeasibility[result_Association,
                   solData_Association] := Module[
  {varNames, sol, freeSyms, allExprs, ineqs},

  varNames = result[``varNames''];
  sol = solData[``solution''];
  freeSyms = solData[``freeSyms''];

  allExprs = Table[sol[vn], {vn, varNames}];
  ineqs = And @@ ((# >= 0) & /@ allExprs);

  If[Length[freeSyms] == 1,
    Reduce[ineqs, freeSyms[[1]], Reals],
    Reduce[ineqs, freeSyms, Reals]]
]
\end{lstlisting}

\subsection{Example invocation (Family A)}

\begin{lstlisting}
result  = solveKL[{"10010``,``11101'',''01100"},
                  {"11000``,``10111'',''00110"}];
solData = extractSolution[result];
zData   = computeZMarginals[result, solData];
region  = computeFeasibility[result, solData];

(* SigmaZ range *)
t1 = solData[``freeSyms''][[1]];
Minimize[{zData[``sigmaZ''], region}, t1]
Maximize[{zData[``sigmaZ''], region}, t1]
\end{lstlisting}

%%%%%%%%%%%%%%%%%%%%%%%%%%%%%%%%%%%%%%%%%%%%%%%%%%%%%%%%%%%%%%%%
\section*{Conclusion}
%%%%%%%%%%%%%%%%%%%%%%%%%%%%%%%%%%%%%%%%%%%%%%%%%%%%%%%%%%%%%%%%

Every algebraic and numerical claim in \texttt{LPv3.tex}---parametric solutions, feasibility regions, $Z$-marginal vectors, $\SigZ$ polynomials and their ranges, and the convex-hull dimension classification---has been independently verified using exact rational arithmetic in the Wolfram Language.
The verification covers all six examples (Families A--E and the worked $((5,2,2))$ example).
Zero discrepancies were found.

\end{document}